\chapter{Different design approaches}
\label{c:approaches}

%%%%%%%%%%%%%%%%%%%%%%%%%%%%%%%%%%%%%%%%%%%%%%%%%%%%%%%%%%%%%%%%%%%%%%%%%%%%%%%%

\section{Approach A: Operator-Algebra}
\begin{outline}
	\1 Concept: Implementation of the exact Algebra
	\1 Structure: Addition, subtraction und multiplication of chained operators.
\end{outline}

This approach tries to follow the exact definition of the algebra operators. It begins with defining the operators itself. A function is defined that encodes the logic of the operators. Given a specific state it calculates how the operator by definition acts on this state. Given the opaerator \(a\dagger\) and the state \(\phi\) it would act as following. \\
\begin{equation}
	a_j^\dagger \phi = a_j^\dagger |n_0, n_1, \dots, n_j, \dots\rangle = (1 - n_j) (-1)^{\sum_{k=0}^{j-1} n_k} |n_0, n_1, \dots, 1_j, \dots\rangle
\end{equation}



\section{Approach B: Generator-based evolution (symbolically)}
\begin{outline}
	\1 Concept: Apply Generator directly on a state 
	\1 Unitary operation: How to calculate $U = \exp(\sum \theta_{pq} \hat{E}_{pq})$ efficient?
\end{outline}